% Emacs, this is -*-latex-*

\ex{Gradiente de fun\'c\~oes vetoriais}
\label{ex:grad-vector}
Para uma fun\'c\~ao $J$ que mapeia um vetor coluna $\w \in \mathbb{R}^n$ para $\mathbb{R}$, o gradiente \'e definido como
\begin{equation}
    \nabla J(\w) = \left(
    \begin{array}{c}
        \frac{\partial J(\w)}{\partial w_1} \\ %
        \vdots \\
        \frac{\partial J(\w)}{\partial w_n}
    \end{array} \right),
\end{equation}
onde $\partial J(\w)/\partial w_i$ s\~ao as derivadas parciais de $J(\w)$ em
rela\'c\~ao ao $i$-\'esimo elemento do vetor $\w = (w_1, \ldots, w_n)^\top$ (na
base can\^onica). Alternativamente, ele \'e definido como sendo o vetor coluna $\nabla J(\w)$ tal que
\begin{align}
    J(\w+\epsilon \h) &=  J(\w) + \epsilon \left( \nabla J(\w)\right)^\top \h+O(\epsilon^2) \label{eq:grad-vector-alt}
\end{align}
para uma perturba\'c\~ao arbitr\'aria $\epsilon \h$. Isso expressa a derivada em
termos de uma aproxima\'c\~ao de primeira ordem, ou afim, para a fun\'c\~ao perturbada
$J(\w+\epsilon \h)$. A derivada $\nabla J$ \'e uma transforma\'c\~ao linear que
mapeia $\h \in \mathbb{R}^n$ para $\mathbb{R}$ \citep[veja, e.g.,][Cap\'itulo 9, para um
tratamento formal de derivadas]{Rudin1976}.

Use qualquer uma das defini\'c\~oes para determinar $\nabla J(\w)$ para as seguintes fun\'c\~oes onde $\a \in \mathbb{R}^n$, $\A \in \mathbb{R}^{n \times n}$ e
$f: \mathbb{R} \rightarrow \mathbb{R}$ \'e uma fun\'c\~ao diferenci\'avel.
\begin{exenumerate}
    \item $J(\w)=\a^\top \w$.
    
    \begin{solution}
        Primeiro m\'etodo:
        \begin{align}
            J(\w) &= \a^\top\w = \sum_{k=1}^n a_kw_k& \implies&& \frac{\partial{J(\w)}}{\partial{w_i}} &= a_i
        \end{align}
        Portanto
        \begin{equation}
            \nabla J(\w) = \begin{pmatrix} a_1 \\ a_2 \\ \vdots \\ a_n \end{pmatrix} = \a.
        \end{equation}
        Segundo m\'etodo:
        \begin{align}
            J(\w + \epsilon\h) &= \a^\top(\w + \epsilon\h) = \underbrace{\a^\top\w}_{J(\w)} + \epsilon \underbrace{\a^\top \h}_{\nabla J^\top\h}
        \end{align}
        Assim, encontramos novamente $\nabla J(\w) = \a. $
    \end{solution}
    
    \item $J(\w)=\w^\top \A\w$.
    
    \begin{solution}
        Primeiro m\'etodo: Come\'camos com
        \begin{align}
            J(\w) &= \w^\top\A\w = \sum_{i=1}^n \sum_{j=1}^n w_iA_{ij}w_j
        \end{align}
        Portanto,
        \begin{align}
            \frac{\partial{J(\w)}}{\partial{w_k}} &= \sum_{j=1}^n A_{kj}w_j + \sum_{i=1}^n w_i A_{ik} \\
            &= \sum_{j=1}^n A_{kj}w_j + \sum_{i=1}^n w_i (\A^\top)_{ki}\\
            &= \sum_{j=1}^m \left(A_{kj}+(\A^\top)_{kj}\right)w_j
        \end{align}
        onde usamos que a entrada na linha $i$ e coluna $k$ da matriz $\A$ \'e igual \`a entrada na linha $k$ e coluna $i$ de sua transposta $\A^\top$. Segue que
        \begin{align}
            \nabla J(\w) &=
            \begin{pmatrix}
                \sum_{j=1}^n  \left(A_{1j}+(\A^\top)_{1j}\right)w_j\\
                \vdots \\
                \sum_{j=1}^n  \left(A_{nj}+(\A^\top)_{nj}\right)w_j
            \end{pmatrix} \\
            &= (\A+\A^\top)\w,
        \end{align}
        onde usamos que somas como $\sum_j B_{ij} w_j$ s\~ao iguais ao $i$-\'esimo elemento do produto matriz-vetor $\B \w$.
        
        Segundo m\'etodo:
        \begin{align}
            J(\w + \epsilon\h) &= (\w + \epsilon\h)^\top \A(\w + \epsilon\h) \\
            &= \w^\top\A\w + \w^\top\A(\epsilon\h) + \epsilon\h^\top\A\w +
            \underbrace{\epsilon\h^\top \A\epsilon\h}_{O(\epsilon^2)} \\
            &= \w^\top\A\w + \epsilon(\w^\top\A\h + \w^\top\A^\top\h) + O(\epsilon^2) \\
            &= \underbrace{\w^\top\A\w}_{J(\w)} +\epsilon(\underbrace{\w^\top\A + \w^\top\A^\top}_{\nabla J(\w)^\top})\h + O(\epsilon^2)
        \end{align}
        onde usamos que $\h^\top\A\w$ \'e um escalar, de modo que $\h^\top\A\w = (\h^\top\A\w)^\top = \w^\top \A^\top \h$. Portanto
        \begin{equation}
            \nabla J(\w)^\top = \w^\top\A + \w^\top\A^\top = \w^\top(\A+\A^\top)
        \end{equation}
        e
        \begin{equation}
            \nabla J(\w) = (\A + \A^\top)\w.
        \end{equation}
    \end{solution}
    
    \item $J(\w)=\w^\top \w$.
    
    \begin{solution}
        A maneira mais f\'acil de calcular o gradiente de $J(\w) = \w^\top\w$ \'e usar a quest\~ao anterior com $\A = \I$ (a matriz identidade).
        Portanto
        \begin{equation}
            \nabla J(\w) = \I\w + \I^\top\w = \w + \w = 2\w.
        \end{equation}
        
    \end{solution}
    
    \item $J(\w)=||\w||_2$.
    
    \begin{solution}
        Note que $||\w||_2 = \sqrt{\w^\top \w}$.
        
        Primeiro m\'etodo: Usamos a regra da cadeia
        \begin{equation}
            \frac{\partial{J(\w)}}{\partial{w_k}} = \frac{\partial \sqrt{\w^\top \w}}{\partial \w^\top \w}\frac{\partial \w^\top \w}{\partial w_k}
        \end{equation}
        e que
        \begin{equation}
            \frac{\partial \sqrt{\w^\top \w}}{\partial \w^\top \w} = \frac{1}{2\sqrt{\w^\top \w}}
        \end{equation}
        As derivadas $\partial \w^\top \w/\partial w_k$ foram calculadas na quest\~ao acima, de modo que
        \begin{equation}
            \nabla J(\w) = \frac{1}{2\sqrt{\w^\top \w}} 2 \w = \frac{\w}{||\w||_2}
        \end{equation}
        Segundo m\'etodo: Seja $f(\w) = \w^\top \w$. Da quest\~ao anterior, sabemos que
        \begin{align}
            f(\w + \epsilon \h) & = f(\w) + \epsilon 2\w^\top \h + O(\epsilon^2).
        \end{align}
        Al\'em disso,
        \begin{align}
            \sqrt{z + \epsilon u + O(\epsilon^2)} & = \sqrt{z} + \frac{1}{2\sqrt{z}}(\epsilon u + O(\epsilon^2)) + O(\epsilon^2)\\
            & = \sqrt{z} + \epsilon \frac{1}{2\sqrt{z}} u + O(\epsilon^2)
        \end{align}
        Com $z=f(\w)$ e $u = 2 \w^\top \h$, obtemos assim
        \begin{align}
            J(\w + \epsilon \h) &= \sqrt{f(\w+\epsilon \h)} \\
            &= \sqrt{f(\w)} + \epsilon \frac{1}{2\sqrt{f(\w)}} 2\w^\top \h + O(\epsilon^2)\\
            &= \sqrt{f(\w)} + \epsilon \frac{\w^\top}{\sqrt{f(\w)}}\h + O(\epsilon^2)\\
            &= J(\w) +  \epsilon \frac{\w^\top}{\sqrt{||\w||_2}}\h + O(\epsilon^2)
        \end{align}
        de modo que
        \begin{equation}
            \nabla J(\w) = \frac{\w}{||\w||_2}.
        \end{equation}
        
    \end{solution}
    
    \item $J(\w)=f(||\w||_2)$.
    \begin{solution}
        Tanto a regra da cadeia quanto a abordagem com a expans\~ao de Taylor podem ser usadas
        para lidar com a fun\'c\~ao externa $f$. Em qualquer caso:
        \begin{align}
            \nabla J(\w) &= f'(||\w||_2) \nabla ||\w||_2 = f'(||\w||_2) \frac{\w}{||\w||_2},
        \end{align}
        onde $f'$ \'e a derivada da fun\'c\~ao $f$.
    \end{solution}
    
    \item $J(\w)=f(\w^\top\a)$.
    
    \begin{solution}
        Vimos que $\nabla_\w \a^\top\w = \a$. O uso da regra da cadeia resulta em
        \begin{align}
            \nabla J(\w)  &= f'(\w^\top\a)\nabla (\w^\top\a) \\
            &= f'(\w^\top\a) \a
        \end{align}
        
        
    \end{solution}
    
\end{exenumerate}

% --------------------------------------------------

\ex{M\'etodo de Newton}
\label{ex:newtons_method}
Assuma que na vizinhan\'ca de $\w_0$, uma fun\'c\~ao $J(\w)$ possa ser descrita pela aproxima\'c\~ao quadr\'atica
\begin{equation}
    f(\w)=c+ \g^\top(\w-\w_0)+\frac{1}{2}(\w-\w_0)^\top \H(\w-\w_0),
\end{equation}
onde $c=J(\w_0)$, $\g$ \'e o gradiente de $J$ em rela\'c\~ao a $\w$, e $\H$ uma
matriz sim\'etrica positiva definida (e.g., a matriz Hessiana de $J(\w)$ em
$\w_0$, se for positiva definida).

\begin{exenumerate}
    \item Use a \exref{ex:grad-vector} para determinar $\nabla f(\w)$.
    
    \begin{solution}
        Primeiro escrevemos $f$ como
        \begin{align}
            f(\w) &= c + \g^\top(\w-\w_0)+\frac{1}{2}(\w-\w_0)^T\H(\w-\w_0) \\
            &= c - \g^\top \w_0 + \frac{1}{2} \w_0^\top\H \w_0 + \nonumber \\
            &\phantom{=} \g^\top \w + \frac{1}{2}\w^\top\H\w - \frac{1}{2} \w_0^\top \H \w - \frac{1}{2} \w^\top \H \w_0
        \end{align}
        Usando agora que $\w^\top \H \w_0$ \'e um escalar e que $\H$ \'e sim\'etrica, temos
        \begin{equation}
            \w^\top \H \w_0= (\w^\top \H \w_0)^\top = \w_0^\top \H^\top \w = \w_0^\top \H \w
        \end{equation}
        e portanto
        \begin{equation}
            f(\w) = \text{const} + (\g^\top - \w_0^\top \H) \w +  \frac{1}{2}\w^\top\H\w
        \end{equation}
        Com os resultados da \exref{ex:grad-vector} e o fato de que $\H$ \'e sim\'etrica, obtemos assim
        \begin{align}
            \nabla f(\w) &= \g-\H^\top \w_0 + \frac{1}{2}(\H^\top\w + \H\w) \\
            & = \g-\H \w_0 + \H\w
        \end{align}
        
        A expans\~ao de $f(\w)$ devido aos termos $\w-\w_0$ \'e um pouco
        tediosa. \'E mais simples notar que os gradientes definem uma aproxima\'c\~ao linear
        da fun\'c\~ao. Podemos lidar de forma mais eficiente com $\w-\w_0$ mudando as
        coordenadas e determinando a aproxima\'c\~ao linear de $f$ como uma fun\'c\~ao de
        $\v = \w-\w_0$, i.e., localmente em torno do ponto $\w_0$. Temos ent\~ao
        \begin{align}
            \tilde{f}(\v) & = f(\v+\w_0)\\
            & = c + \g^\top \v + \frac{1}{2} \v^\top \H \v
        \end{align}
        Com a \exref{ex:grad-vector}, a derivada \'e
        \begin{equation}
            \nabla_{\v} \tilde{f}(\v) = \g + \H \v
        \end{equation}
        e a aproxima\'c\~ao linear torna-se
        \begin{equation}
            \tilde{f}(\v+\epsilon \h) =  c + \epsilon (\g + \H \v)^\top \h + O(\epsilon^2)
        \end{equation}
        A aproxima\'c\~ao linear para $\tilde{f}$ determina uma aproxima\'c\~ao linear de $f$ em torno de $\w_0$, i.e.,
        \begin{equation}
            f(\w + \epsilon \h) =  \tilde{f}(\w-\w_0+\epsilon \h) =  c + \epsilon (\g + \H (\w-\w_0))^\top \h + O(\epsilon^2)
        \end{equation}
        de modo que a derivada para $f$ \'e
        \begin{equation}
            \nabla_{\w}f(\w) = \g + \H (\w-\w_0) = \g - \H \w_0 + \H \w,
        \end{equation}
        que \'e o mesmo resultado de antes.
    \end{solution}
    
    \item Uma condi\'c\~ao necess\'aria para que $\w$ seja \'otimo (levando a um m\'aximo, m\'inimo ou ponto de sela) \'e $\nabla f(\w)=0$. 
    Determine $\w^\ast$ tal que $\nabla f(\w)\big|_{\w=\w^\ast}=0$. Forne\'ca argumentos do porqu\^e $\w^\ast$ \'e um minimizador de $f(\w)$.
    
    \begin{solution}
        Definimos o gradiente como zero e resolvemos para $\w$:
        \begin{equation}
            \g + \H(\w -\w_0) = 0 \quad \leftrightarrow \quad \w-\w_0 = - \H^{-1} \g
        \end{equation}
        de modo que
        \begin{equation}
            \w^\ast = \w_0 - \H^{-1} \g.
        \end{equation}
        Como assumimos que $\H$ \'e positiva definida, a inversa de $\H$ existe (e
        tamb\'em \'e positiva definida).
        
        Vamos considerar $f$ como uma fun\'c\~ao de $\v$ em torno de $\w^\ast$, i.e., $\w = \w^\ast+\v$. Com $\w^\ast+\v-\w_0= -\H^{-1}\g + \v$, temos
        \begin{align}
            f(\w^\ast+\v) & = c + \g^\top(- \H^{-1} \g+\v) + \frac{1}{2}(-\H^{-1}\g + \v)^\top \H (-\H^{-1}\g + \v)
        \end{align}
        Como $\H$ \'e positiva definida, temos que $(-\H^{-1}\g + \v)^\top \H
        (-\H^{-1}\g + \v)>0$ para todo $\v$. Assim, \`a medida que nos afastamos de $\w^\ast$,
        a fun\'c\~ao aumenta quadraticamente, de modo que $\w^\ast$ minimiza $f(\w)$.
        
    \end{solution}
    
    
    \item Em termos do m\'etodo de Newton para minimizar $J(\w)$, o que $\w_0$ e $\w^\ast$ representam?
    
    \begin{solution}
        A equa\'c\~ao
        \begin{equation}
            \w^\ast = \w_0 - \H^{-1} \g.
        \end{equation}
        corresponde a um passo de atualiza\'c\~ao no m\'etodo de Newton, onde $\w_0$ \'e o
        valor atual de $\w$ na otimiza\'c\~ao de $J(\w)$ e $\w^\ast$ \'e o
        valor atualizado. Na pr\'atica, em vez de determinar a inversa $\H^{-1}$, resolvemos
        \begin{equation}
            \H \p = \g
        \end{equation}
        para $\p$ e ent\~ao definimos $\w^\ast = \w_0 - \p$. O vetor $\p$ \'e a
        dire\'c\~ao de busca, e \'e poss\'ivel incluir um tamanho de passo $\alpha$ para que
        a atualiza\'c\~ao se torne $\w^\ast = \w_0 - \alpha \p$. O valor de $\alpha$ pode
        ser definido manualmente ou determinado via m\'etodos de busca linear \citep[veja
        e.g.,][]{Nocedal1999}.
    \end{solution}
\end{exenumerate}

% --------------------------------------------------

\ex{Gradiente de fun\'c\~oes matriciais}
\label{ex:grad-matrix}
Para fun\'c\~oes $J$ que mapeiam uma matriz $\W \in \mathbb{R}^{n\times m}$ para $\mathbb{R}$, o gradiente \'e definido como 
\begin{equation}
    \nabla J(\W) = 
    \begin{pmatrix}
        \frac{\partial J(\W)}{\partial W_{11}} & \ldots &  \frac{\partial J(\W)}{\partial W_{1m}} \\
        \vdots & \vdots & \vdots\\
        \frac{\partial J(\W)}{\partial W_{n1}} & \ldots &  \frac{\partial J(\W)}{\partial W_{nm}}
    \end{pmatrix}.
\end{equation}

Alternativamente, ele \'e definido como sendo a matriz $\nabla J$ tal que
\begin{align}
    J(\W+\epsilon \H) &= J(\W) + \epsilon \tr(\nabla J^\top \H) + O(\epsilon^2) \label{eq:grad-matrix-alt}\\
    &= J(\W) + \epsilon \tr(\nabla J \H^\top) + O(\epsilon^2)
\end{align}
Essa defini\'c\~ao \'e an\'aloga \`a das fun\'c\~oes vetoriais
em \eqref{eq:grad-vector}. Ela expressa a derivada em termos de uma aproxima\'c\~ao
linear para o objetivo perturbado $J(\W+\epsilon \H)$ e, mais formalmente,
$\tr \nabla J^\top$ \'e uma transforma\'c\~ao linear que mapeia
$\H \in \mathbb{R}^{n \times m}$ para $\mathbb{R}$ \citep[veja e.g.,][Cap\'itulo 9, para um tratamento formal de derivadas]{Rudin1976}.

Seja $\e^{(i)}$ um vetor \emph{coluna} que \'e zero em todos os lugares, exceto na posi\'c\~ao $i$
onde vale 1. Al\'em disso, seja $\e^{[j]}$ um vetor \emph{linha} que \'e
zero em todos os lugares, exceto na posi\'c\~ao $j$ onde vale 1. O produto externo
$\e^{(i)}\e^{[j]}$ \'e ent\~ao uma matriz que \'e zero em todos os lugares, exceto na linha $i$ e
coluna $j$ onde vale um. Para $\H = \e^{(i)} \e^{[j]}$, obtemos
\begin{align}
    J(\W+\epsilon \e^{(i)}\e^{[j]}) &= J(\W) + \epsilon \tr((\nabla J)^\top \e^{(i)} \e^{[j]}) + O(\epsilon^2)\\
    &= J(\W) + \epsilon  \e^{[j]} (\nabla J)^\top \e^{(i)} + O(\epsilon^2)\\
    &= J(\W) + \epsilon \e^{[i]} \nabla J  \e^{(j)}+O(\epsilon^2)
\end{align}
Note que $\e^{[i]} \nabla J \e^{(j)}$ seleciona o elemento da
matriz $\nabla J$ que est\'a na linha $i$ e coluna $j$, i.e., $\e^{[i]} \nabla J
\e^{(j)}=\partial J/\partial W_{ij}$.

Use qualquer uma das duas defini\'c\~oes para encontrar $\nabla
J(\W)$ para as fun\'c\~oes abaixo, onde
$\u \in \mathbb{R}^n, \v \in \mathbb{R}^m, \A \in \mathbb{R}^{n \times m}$, e
$f: \mathbb{R} \rightarrow \mathbb{R}$ \'e diferenci\'avel.

\begin{exenumerate}
    \item $J(\W) = \u^\top\W\v$.
    
    \begin{solution}
        Primeiro m\'etodo: Com $J(\W) = \sum_{i=1}^n \sum_{j=1}^m u_iW_{ij}v_j$ temos
        \begin{align}
            \frac{\partial J(\W)}{W_{kl}} &= u_kv_l = (\u\v^\top)_{kl}
        \end{align}
        e portanto
        \begin{equation}
            \nabla J(\W) = \u\v^\top
        \end{equation}
        
        Segundo m\'etodo: 
        \begin{align}
            J(\W + \epsilon \H) & = \u^\top (\W+\epsilon \H)\v\\
            & = J(\W) + \epsilon \u^\top \H \v\\
            & = J(\W) + \epsilon \tr(\u^\top \H \v)\\
            & = J(\W) + \epsilon \tr(\v \u^\top \H)
        \end{align}
        Portanto:
        \begin{align}
            \nabla J(\W) & = \u\v^\top
        \end{align}
        
    \end{solution}
    
    \item $J(\W) = \u^\top(\W+\A)\v$.
    \begin{solution}
        A expans\~ao da fun\'c\~ao objetivo resulta em $J(\W) = \u^\top \W \v
        + \u^\top \A \v$. O segundo termo n\~ao depende de $\W$. Pela quest\~ao anterior,
        a derivada, portanto, \'e
        \begin{equation}
            \nabla J(\W) = \u \v^\top
        \end{equation}
        
    \end{solution}
    \item $J(\W) = \sum_n f(\w_n^\top\v)$, onde $\w_n^\top$ s\~ao as linhas da matriz $\W$.
    
    \begin{solution}
        Primeiro m\'etodo:
        \begin{align}
            \frac{\partial J(\W)}{\partial W_{ij}} &= \sum_{k=1}^n \frac{\partial}{\partial W_{ij}} f(\w_k^\top\v) \\
            &= f'(\w_i^\top\v) \frac{\partial}{\partial W_{ij}} \underbrace{\w_i^\top\v}_\text{$\sum_{j=1}^m W_{ij}v_j$} \\
            &= f'(\w_i^\top\v)v_j
        \end{align}
        Portanto
        \begin{equation}
            \nabla J(\W) = f'(\W\v)\v^\top,
        \end{equation}
        onde $f'$ opera elemento a elemento no vetor $\W\v$.
        
        Segundo m\'etodo:
        \begin{align}
            J(\W) &= \sum_{k=1}^n f(\w_k^\top\v)\\
            &= \sum_{k=1}^n f(\e^{[k]}\W\v),
        \end{align}
        onde $\e^{[k]}$ \'e o vetor linha unit\'ario que \'e zero em todos os lugares, exceto pelo elemento $k$ que \'e igual a um. Agora realizamos uma perturba\'c\~ao de $\W$ por $\epsilon \H$.
        \begin{align}
            J(\W + \epsilon \H) &= \sum_{k=1}^n f(\e^{[k]}(\W + \epsilon \H)\v) \\
            &= \sum_{k=1}^n f(\e^{[k]}\W\v + \epsilon\e^{[k]}\H\v) \\
            &= \sum_{k=1}^n (f(\e^{[k]}\W\v) + \epsilon f'(\e^{[k]}\W\v)\e^{[k]}\H\v + O(\epsilon^2)\\
            &= J(\W) + \epsilon \left( \sum_{k=1}^n  f'(\e^{[k]}\W\v)\e^{[k]}\right) \H \v + O(\epsilon^2)
        \end{align}
        O termo $f'(\e^{[k]}\W\v) \e^{[k]}$ \'e um vetor linha que \'e igual a $(0, \ldots, 0,f'(\e^{[k]}\W\v), 0, \ldots, 0)$. Assim, temos
        \begin{equation}
            \sum_{k=1}^n  f'(\e^{[k]}\W\v)\e^{[k]}) = f'(\W\v)^\top
        \end{equation}
        onde $f'$ opera elemento a elemento no vetor coluna $\W\v$. A fun\'c\~ao objetivo perturbada \'e, portanto
        \begin{align}
            J(\W + \epsilon \H) & = J(\W) + \epsilon f'(\W\v)^\top \H \v  + O(\epsilon^2)\\
            & =  J(\W) + \epsilon \tr\left(f'(\W\v)^\top \H \v\right)  + O(\epsilon^2)\\
            & =  J(\W) + \epsilon \tr\left(\v f'(\W\v)^\top \H \right)  + O(\epsilon^2)
        \end{align}
        Portanto, o gradiente \'e a transposta de $\v f'(\W\v)^\top$, i.e.,
        \begin{equation}
            \nabla J(\W) = f'(\W\v) \v^\top
        \end{equation}
    \end{solution}
    
    \item $J(\W) = \u^\top \W^{-1}\v$ (\emph{Dica: $(\W+\epsilon \H)^{-1} = \W^{-1} -\epsilon \W^{-1}\H\W^{-1}+O(\epsilon^2)$.})
    \begin{solution}
        Primeiro verificamos a dica:
        \begin{align}
            \left(\W^{-1} -\epsilon \W^{-1}\H\W^{-1}+O(\epsilon^2)\right) (\W+\epsilon \H)&=  \I + \epsilon \W^{-1} \H - \epsilon \W^{-1}\H + O(\epsilon^2)\\
            &= \I + O(\epsilon^2)
        \end{align}
        Assim, a identidade se mant\'em para termos menores que $\epsilon^2$, o que \'e
        suficiente, pois n\~ao nos importamos com termos de ordem $\epsilon^2$ e menores 
        na defini\'c\~ao do gradiente em \eqref{eq:grad-matrix}.
        
        Vamos, portanto, fazer uma aproxima\'c\~ao de primeira ordem do objetivo perturbado $J(\W + \epsilon \H)$: 
        \begin{align}
            J(\W + \epsilon\H) &= \u^\top(\W + \epsilon\H)^{-1}\v \\
            &\overset{\textrm{dica}}{=} \u^\top(\W^{-1} - \epsilon \W^{-1}\H \W^{-1} + O(\epsilon^2))\v \\
            &= \u^\top \W^{-1}\v - \epsilon \u^\top\W^{-1}\H\W^{-1}\v + O(\epsilon^2) \\
            &= J(\W) -  \epsilon \tr\left(\u^\top\W^{- 1}\H\W^{-1}\v\right) + O(\epsilon^2) \\
            &= J(\W) - \epsilon \tr\left(\W^{-1}\v\u^\top\W^{-1}\H\right) + O(\epsilon^2)
        \end{align}
        A compara\'c\~ao com \eqref{eq:grad-matrix} fornece
        \begin{equation}
            \nabla J^\top = -\W^{-1}\v\u^\top\W^{-1}
        \end{equation}
        e portanto
        \begin{equation}
            \nabla J = -\W^{-\top}\u\v^\top\W^{-\top},
        \end{equation}
        onde $\W^{-\top}$ \'e a transposta da inversa de $\W$.
        
    \end{solution}
\end{exenumerate}

% --------------------------------------------------

\ex{Gradiente do log-determinante}
\label{ex:grad-log-det}
O objetivo deste exerc\'icio \'e determinar o gradiente de
\begin{equation}
    J(\W) = \log | \det(\W)|.
\end{equation}

\begin{exenumerate}
    \item Mostre que o $n$-\'esimo autovalor $\lambda_n$ pode ser escrito como
    \begin{equation}
        \lambda_n=\v_n^\top \W \u_n,
    \end{equation}
    onde $\u_n$ \'e o $n$-\'esimo autovetor e $\v_n$ o $n$-\'esimo vetor coluna de $\U^{-1}$, sendo $\U$ a matriz com os autovetores $\u_n$ como colunas.
    
    \begin{solution}
        Como na \exref{ex:eigenvalue-decomposition}, seja $\U\Lambdab \V^\top$ a
        decomposi\'c\~ao em autovalores de $\W$ (com $\V^\top = \U^{-1}$). Ent\~ao $\Lambdab = \V^\top
        \W\U$ e
        \begin{align}
            \lambda_n &= \e^{[n]}\Lambdab\e^{(n)} \\
            &= \e^{[n]}\V^\top \W \U\e^{(n)} \\
            &= (\V\e^{(n)})^\top \W\U\e^{(n)} \\
            &= \v_n^\top \W\u_n,
        \end{align}
        onde $\e^{(n)}$ \'e o vetor da base can\^onica (unit\'ario) com um 1 na $n$-\'esima
        posi\'c\~ao e zeros nos outros lugares, e $\e^{[n]}$ \'e o vetor linha correspondente.
    \end{solution}
    
    \item Calcule o gradiente de $\lambda_n$ em rela\'c\~ao a $\W$, i.e., $\nabla \lambda_n(\W)$.
    
    \begin{solution}
        Pela \exref{ex:grad-matrix}, temos
        \begin{align}
            \nabla_{\W} \lambda_n(\W) &= \nabla_{\W} \v_n^\top \W \u_n = \v_n \u_n^\top.
        \end{align}
    \end{solution}
    
    \item Escreva $J(\W)$ em termos dos autovalores $\lambda_n$ e calcule $\nabla J(\mathbf{\W})$.
    \begin{solution}
        Na \exref{ex:trace-determinants-eigenvalues}, mostramos que $\det(\W) = \prod_i \lambda_i$ e, portanto, $|\textrm{det}(\W)| = \prod_i |\lambda_i|.$
        \begin{enumerate}
            \item[(i)] Se $\W$ for positiva definida, seus autovalores s\~ao positivos e
            podemos omitir os valores absolutos, de modo que $|\det(\W)| = \prod_i \lambda_i$.
            \item[(ii)] Se $\W$ for uma matriz com entradas reais, ent\~ao $\W\u = \lambda\u$
            implica em $\W\bar{\u} = \bar{\lambda}\bar{\u}$, i.e., se $\lambda$ for um
            autovalor complexo, ent\~ao $\bar{\lambda}$ (o conjugado complexo de
            $\lambda$) tamb\'em \'e um autovalor. Como $|\lambda|^2 =
            \lambda\bar{\lambda}$,
            \begin{align}
                |\det(\W)| = \left(\prod_{\lambda_i \in \mathbb{C}} \lambda_i \right) \left(\prod_{\lambda_j \in \mathbb{R}} |\lambda_j| \right).
            \end{align}
        \end{enumerate}
        Agora podemos escrever $J(\W)$ em termos dos autovalores:
        \begin{align}
            J(\W)  &= \log |\det(\W)| \\
            &= \log\left(\prod_{\lambda_i \in \mathbb{C}} \lambda_i \right) \left(\prod_{\lambda_j \in \mathbb{R}} |\lambda_j|\right) \\
            &= \log\left(\prod_{\lambda_i \in \mathbb{C}} \lambda_i \right) + \log\left(\prod_{\lambda_j \in \mathbb{R}} |\lambda_j| \right) \\
            &= \sum_{\lambda_i \in \mathbb{C}} \log \lambda_i + \sum_{\lambda_j \in \mathbb{R}} \log |\lambda_j|.
        \end{align}
        Assuma que os $\lambda_j$ reais s\~ao n\~ao nulos, de modo que
        \begin{align}
            \nabla_{\W} \log |\lambda_j| & = \frac{1}{|\lambda_j|} \nabla_{\W} |\lambda_j|\\
            & =  \frac{1}{|\lambda_j|} \textrm{sign}(\lambda_j)\nabla_{\W}\lambda_j
        \end{align}
        Portanto
        \begin{align}
            \nabla J(\W) &= \sum_{\lambda_i \in \mathbb{C}}  \nabla_{\W}\log \lambda_i + \sum_{\lambda_j \in \mathbb{R}}  \nabla_{\W} \log |\lambda_j|\\ 
            &= \sum_{\lambda_i \in \mathbb{C}} \frac{1}{\lambda_i} \nabla_{\W} \lambda_i + 
            \sum_{\lambda_i \in \mathbb{R}} \frac{1}{|\lambda_i|} \textrm{sign}(\lambda_i)\nabla_{\W}\lambda_i\\
            &= \sum_{\lambda_i \in \mathbb{C}} \frac{\v_i\u_i^\top}{\lambda_i} + 
            \sum_{\lambda_i \in \mathbb{R}} \frac{\textrm{sign}(\lambda_i)\v_i\u_i^\top}{|\lambda_i|} \\
            &= \sum_{\lambda_i \in \mathbb{C}} \frac{\v_i\u_i^\top}{\lambda_i} + 
            \sum_{\lambda_i \in \mathbb{R}} \frac{\v_i\u_i^\top}{\lambda_i} \\
            &= \sum_i  \frac{\v_i\u_i^\top}{\lambda_i}.
        \end{align}
        
    \end{solution}
    \item Mostre que
    \begin{equation}
        \nabla J(\W) = (\W^{-1})^\top.
    \end{equation}
    
    \begin{solution}
        Isso segue da \exref{ex:eigenvalue-decomposition}, onde descobrimos
        que
        \begin{equation}
            \W^{-1}= \sum_i \frac{1}{\lambda_i}\u_i \v_i^\top.
        \end{equation}
        De fato:
        \begin{align}
            \nabla J(\W) = \sum_i  \frac{\v_i\u_i^\top}{\lambda_i} = \sum_i \frac{1}{\lambda_i} (\u_i\v_i^\top)^\top =  \left(\W^{-1}\right)^\top.
        \end{align}
        
    \end{solution}
\end{exenumerate}

% --------------------------------------------------

\ex{Dire\'c\~oes de descida para fun\'c\~oes matriciais}
\label{ex:descent-directions-for-matrix-valued-functions}
Assuma que gostar\'iamos de minimizar uma fun\'c\~ao matricial $J(\W)$ por gradiente
descendente, i.e., a equa\'c\~ao de atualiza\'c\~ao \'e
\begin{equation}
    \W \leftarrow \W - \epsilon \nabla J(\W),
\end{equation}
onde $\epsilon$ \'e o tamanho do passo. O gradiente $\nabla J(\W)$ foi definido na
\exref{ex:grad-matrix}. Foi apontado l\'a que o gradiente define uma
aproxima\'c\~ao de primeira ordem para a fun\'c\~ao objetivo perturbada $J(\W+\epsilon
\H)$. Pela \eqref{eq:grad-matrix},
\begin{align}
    J(\W - \epsilon \nabla J(\W)) & = J(\W) - \epsilon \tr(\nabla J(\W)^\top \nabla J(\W)) + O(\epsilon^2)
\end{align}
Para qualquer matriz $\M$ (n\~ao nula), vale que
\begin{align}
    \tr(\M^\top \M) &= \sum_{i} (\M^\top\M)_{ii}\\
    &= \sum_i \sum_j (\M^\top)_{ij} (\M)_{ji}\\
    &= \sum_i \sum_j M_{ji} M_{ji}\\
    &= \sum_{ij} (M_{ji})^2\\
    &> 0, \label{eq:trace-nonnegativity-alt}
\end{align}
o que significa que $\tr(\nabla J(\W)^\top \nabla J(\W)) > 0$ se o gradiente for n\~ao nulo, e portanto
\begin{equation}
    J(\W - \epsilon \nabla J(\W)) < J(\W)
\end{equation}
para $\epsilon$ suficientemente pequeno. Consequentemente, $\nabla J(\W)$ \'e uma dire\'c\~ao de descida.
Mostre que $\A^\top \A \nabla J(\W)\B\B^\top$ para matrizes n\~ao nulas
$\A$ e $\B$ tamb\'em \'e uma dire\'c\~ao de descida ou deixa o objetivo invariante.

\begin{solution}
    Como na introdu\'c\~ao da quest\~ao, recorremos \`a \eqref{eq:grad-matrix} para
    obter
    \begin{align}
        J(\W - \epsilon \A^\top \A \nabla J(\W) \B\B^\top) & = J(\W) - \epsilon \tr(\nabla J(\W)^\top \A^\top \A \nabla J(\W)\B\B^\top) + O(\epsilon^2)\\
        & =  J(\W) - \epsilon \tr(\B^\top \nabla J(\W)^\top \A^\top \A \nabla J(\W)\B) + O(\epsilon^2),
    \end{align}
    onde $\tr(\B^\top \nabla J(\W)^\top\A^\top \A \nabla J(\W)\B)$ assume a forma
    $\tr(\M^\top \M)$ com $\M= \A \nabla J(\W)\B$. Pela
    \eqref{eq:trace-nonnegativity-alt}, temos assim que $\tr(\B^\top \nabla J(\W)^\top \A^\top
    \A \nabla J(\W)\B) > 0$ se $\A \nabla J(\W)\B$ for n\~ao nulo, e portanto
    \begin{equation}
        J(\W - \epsilon \A^\top \A \nabla J(\W) \B\B^\top) < J(\W)
    \end{equation}
    para $\epsilon$ suficientemente pequeno. Temos igualdade se $\A \nabla J(\W)\B = 0$,
    e.g., se todas as colunas de $\B$ estiverem no espa\'co nulo de $\nabla J$.
\end{solution}